% ============================================================================
%  二维 N=2 理论的相
%  Phases of N=2 Theories in Two Dimensions
%  Edward Witten (School of Natural Sciences, Institute for Advanced Study)
%  IASSNS-HEP-93/3, arXiv:hep-th/9301042v3 [hep-th], 92 pp.
%
%  中文学术译注本：语义级全译 + 公式逐式核对 + 学术审校批注
%
%  编译：xelatex main.tex   （连续两遍，以生成目录与交叉引用）
%  亦可用 tectonic： tectonic -X compile main.tex
%  中文支持使用 ctex 的 fandol 字库（跨平台、无需系统字体）；
%  在 Windows 上若希望改用系统宋体/黑体，将下行改为 fontset=windows。
% ============================================================================

% !TEX program = xelatex
\documentclass[11pt,a4paper]{article}

\usepackage[UTF8,scheme=plain,fontset=fandol]{ctex}

% ---------------------------------------------------------------------------
% 数学与排版
% ---------------------------------------------------------------------------
\usepackage{amsmath}
\usepackage{amssymb}
\usepackage{amsthm}
\usepackage{mathtools}
\usepackage{bm}
\usepackage{mathrsfs}
\usepackage[a4paper,top=2.4cm,bottom=2.4cm,left=2.3cm,right=2.3cm]{geometry}
\usepackage{graphicx}
\usepackage{xcolor}
\usepackage{enumitem}
\usepackage{booktabs}
\usepackage{array}
\usepackage{caption}
\usepackage{fancyhdr}
\usepackage[hidelinks,bookmarksnumbered=true,bookmarksopen=true]{hyperref}

\hypersetup{
    hypertexnames=false,   % 显式 \tag 编号与自动计数器易产生重名锚点，故关闭
    pdftitle={二维 N=2 理论的相（中文译注本）},
    pdfauthor={Edward Witten（原著）；中文译注},
    pdfsubject={arXiv:hep-th/9301042v3, IASSNS-HEP-93/3},
    pdfkeywords={N=2 supersymmetry, linear sigma model, Calabi-Yau,
                 Landau-Ginzburg, toric variety, (0,2) model}
}

% ---------------------------------------------------------------------------
% 版式细节
% ---------------------------------------------------------------------------
\allowdisplaybreaks
\renewcommand{\contentsname}{目\quad 录}
\renewcommand{\abstractname}{摘\quad 要}
\renewcommand{\refname}{参考文献}
\renewcommand{\figurename}{图}
\renewcommand{\tablename}{表}
\setlength{\emergencystretch}{2em}
\setlength{\parskip}{0.2em}

% 目录：一级条目也加引导点（与二级条目风格统一）
\makeatletter
\renewcommand*\l@section{\@dottedtocline{1}{0em}{1.6em}}
\makeatother

% 常用记号（SPEC §4.5 允许的固定写法）
\newcommand{\Tr}{\operatorname{Tr}}
\newcommand{\tr}{\operatorname{tr}}
\newcommand{\dd}{\mathrm{d}}
\newcommand{\ii}{\mathrm{i}}
\newcommand{\ee}{\mathrm{e}}

% 「本节重点」区块：上下细线 + 缩进，可跨页（不依赖额外宏包）
\newenvironment{keytakeaway}
  {\par\medskip\begin{quote}\small\noindent
   \rule{\linewidth}{0.6pt}\par\nopagebreak\smallskip\noindent
   \textbf{本节重点}\par\nobreak\smallskip\noindent\ignorespaces}
  {\par\smallskip\noindent\rule{\linewidth}{0.6pt}\end{quote}\medskip}

% 「译者说明」区块
\newenvironment{translatornote}
  {\begin{quote}\small\noindent\ignorespaces}
  {\end{quote}\medskip}

\pagestyle{fancy}
\fancyhf{}
\fancyhead[L]{\small 二维 $N=2$ 理论的相}
\fancyhead[R]{\small Witten, 1993（中文译注本）}
\fancyfoot[C]{\thepage}
\renewcommand{\headrulewidth}{0.4pt}

% \input 容错：分片尚未交付时不中断编译
\newcommand{\frag}[1]{%
  \IfFileExists{sections/#1.tex}{\input{sections/#1.tex}}{%
    \par\noindent\textcolor{red}{[缺分片：#1]}\par}}

% ---------------------------------------------------------------------------
\begin{document}

\title{\bfseries 二维 $N=2$ 理论的相\\[0.6em]
       \large Phases of $N=2$ Theories in Two Dimensions}
\author{Edward Witten\thanks{原文脚注：本研究部分得到 NSF 基金 PHY92-45317 的资助。}\\[0.3em]
        \small School of Natural Sciences, Institute for Advanced Study\\
        \small Olden Lane, Princeton, NJ 08540}
\date{\small IASSNS-HEP-93/3\quad 1993 年 1 月\\[0.2em]
      \small arXiv:hep-th/9301042v3（1993 年 2 月 18 日）}
\maketitle

\begin{abstract}
\noindent
通过考察超对称规范场论中当参数被改变时所发生的相变，我们找到了建立在加权射影空间中的
Calabi--Yau 超曲面上的 $\sigma$ 模型与 Landau--Ginzburg 模型之间的一种自然联系。
这一构造使人能够重新导出这两类模型之间已知的对应关系，并把它大幅推广，以纳入新的流形类别，
也纳入具有 $(0,2)$ 世界面超对称的模型。该构造还预言了某些涉及时空拓扑改变的物理过程是可能的。
\end{abstract}

\begin{translatornote}
\textbf{译者说明。}
本译本的唯一目标是「无损」：正文逐句对译，公式逐式对照原文页面图像核对，
一切编号（公式 $(x.y)$、文献 $[n]$、章节号）均保持原样。
原文脚注在译文中保留为脚注，其编号由排版系统重新生成，故与原文页码不同，此为正常现象。
文中以下列标记形式插入审校批注（均为 \LaTeX{} 注释，不影响排版输出，可在源码中直接阅读）：

\begin{itemize}[leftmargin=2em,itemsep=1pt,topsep=2pt]
\item \texttt{\% [Tutor Note: \dots]} —— 推导补足、易错预警或物理图像说明；
\item \texttt{\% [原文疑误: \dots]} —— 已判定为原文笔误处，附判断依据与建议；
\item \texttt{\% [存疑: \dots]} —— 无法从图像确证、但疑有问题之处；
\item \texttt{\% [背景补充: \dots]} —— 对中文读者构成实质门槛的术语或背景一句话说明。
\end{itemize}

原文（arXiv:hep-th/9301042v3）全篇 92 页，经矢量图形与位图检测核实为
\textbf{无插图、无带框表格}的纯文字论文，因此本译本不含任何从 PDF 剥离的图像文件；
图、表在原文本就不存在，而非翻译时遗漏。文献列表保留英文原貌（书目信息不可翻译）。
\end{translatornote}

\clearpage
\tableofcontents
\clearpage

% ---------------------------------------------------------------------------
% 正文分片
% ---------------------------------------------------------------------------
\frag{s01-intro}
\frag{s02a-fieldtheory}
\frag{s02b-fieldtheory}
\frag{s03a-model}\frag{s03b-singularity}   % 两片衔接处为句中断句，必须同行（s03a 行末带 %）
\frag{s03c-applications}
\frag{s03d-twisted}
\frag{s04-geometry}
\frag{s05a-wps-toric}\frag{s05a2-comment}   % p58 顶部 §5.2 收尾 + "A General Comment"（句中断句，同行）
\frag{s05b-grassmann-intersect}
\frag{s05c-topology}
\frag{s06a-superfields}
\frag{s06b-models}
\frag{s07-references}

% 附录：原文页码—译文页码对照（由 tools/gen_pagemap.py 生成）
\frag{appendix-pagemap}

\end{document}
